\documentclass[10pt,a4paper]{amsart}
\usepackage[margin=19mm]{geometry}
\usepackage{amsmath,amssymb,mathtools}
\usepackage[scr=rsfs,cal=euler]{mathalfa}
\usepackage{xfrac}
\usepackage[unicode,hidelinks,bookmarks=false]{hyperref}
\newcommand{\ie}{i.e.\ }
\newcommand{\cf}{cf.\ }
\newcommand{\resp}{respectively\ }
\newcommand{\Subsection}{\subsection}
\providecommand{\nobreakdash}{\nobreak\hbox{-}}
\setlength{\emergencystretch}{2em}
\multlinegap0pt
\newcommand{\Spec}{\operatorname{Spec}}
\def\spzb{\overline{\Spec \Z}}
\usepackage{tikz}
\usetikzlibrary{cd}
\usepackage[shortlabels]{enumitem}
\usepackage{mathtools}
\usepackage{amssymb}
\usepackage{float}
\newtheorem{lem}{Lemma}
\newcommand{\A}{\mathbb{A}}
\def\Hom{\rm{Hom}}
\newcommand{\Q}{\mathbf{Q}}
\newcommand{\Z}{\mathbf{Z}}
\newcommand{\Zhat}{\widehat{\Z}}
\title{On the Jacobian of $\spzb$}
\author{Alain Connes, Caterina Consani}
\date{Source: arXiv:2602.15941v1 (CC BY 4.0)}
\begin{document}
\maketitle

\noindent\emph{Source and licence.} On the Jacobian of $\spzb$. Version arXiv:2602.15941v1 (CC BY 4.0), distributed by arXiv under the Creative Commons Attribution 4.0 International licence, http://creativecommons.org/licenses/by/4.0/, as recorded in the arXiv licence metadata for this exact version and shown on the abstract page https://arxiv.org/abs/2602.15941v1. Full licence notice: \texttt{licenses/arxiv-2602.15941-CC-BY-4.0.txt}.\par\smallskip

\noindent\textit{Editorial scope.} Counted unit: the article's lem isopsi (source0.tex lines 2648--2650). The article's complete original proof of it is delivered with it (source0.tex lines 2652--2715, one \texttt{proof} environment of 64 lines). No mathematics is written, repaired or re-derived: the statement and the proof are the article's own lines, in the article's own notation, and no retained line is substituted or re-flowed.  Retained uncounted as the source-local context the proof needs: lines 2644--2647.  Nothing was omitted from any retained span, so the package carries no omission record.  No retained line contains a citation, so the wrapper carries no bibliography and no bibliography entry.  The retained lines contain the article's own diagram code, which the wrapper typesets with the standard tikz packages; no image file is delivered and the package declares no asset dependency.  Editorial formatting only, in addition: an AMS \texttt{amsart} preamble that restates the article's own declarations and macros used by the retained lines, namely the environments \texttt{lem} on separate counters, together with the standard packages those lines need.  The retained environments print in this document as Lemma 1 in order of appearance, whereas the article prints the same items with the numbering of its own full document.  The complete native source of the article as distributed in the arXiv e-print for this exact version is bundled unchanged as \texttt{sources/194837/source0.tex}.
\begin{equation} \label{quotad1}
\psi: \A_f \to \Hom(\Q, \Q/\Z)\quad 
\psi(a)(q) = \pi(qa) \quad \text{for } a \in \A_f,~ q \in \Q.
\end{equation}

\begin{lem}\label{isopsi}
The map $\psi$ is an isomorphism of topological groups.
\end{lem}

\begin{proof}
We proceed by establishing injectivity and surjectivity using the structural properties of $\A_f$ and $\Q$.

\emph{Injectivity.} 
Suppose $\psi(a) = 0$ for some $a \in \A_f$. By definition, this means:
$
\pi(qa) = 0~ \forall q \in \Q
$. 
This implies that $qa \in \ker(\pi) = \Zhat~ \forall q \in \Q$. The following cases arise:
\begin{itemize}[leftmargin=*, labelindent=0pt]
    \item[] If $q=1$, then  $a \in \Zhat$.
    \item[] If $q=1/n$ for an integer $n$, then $a/n \in \Zhat$, which implies $a \in n\Zhat$.
\end{itemize}
Thus, we conclude that $a \in \bigcap_{n \in \Z \setminus \{0\}} n\Zhat$.
Since $\Zhat \cong \prod_p \Z_p$ and $\bigcap_k p^k \Z_p = \{0\}$, it follows that $a = 0$.

\emph{Surjectivity.} 
We compare the defining exact sequence of $\A_f$ 
\begin{equation}\label{seqA}
0 \longrightarrow \Zhat \longrightarrow \A_f \xrightarrow{\pi} \Q/\Z \longrightarrow 0
\end{equation}
with the exact sequence obtained by applying the functor $\Hom(-, \Q/\Z)$ to the defining sequence of $\Q$ 
\[0 \to \Z \to \Q \to \Q/\Z \to 0.
\]
Since $\Q/\Z$ is divisible (injective in the category of abelian groups), the functor $\Hom(-, \Q/\Z)$ is exact, thus we have
\begin{equation}\label{seqB}
0 \longrightarrow \Hom(\Q/\Z, \Q/\Z) \longrightarrow \Hom(\Q, \Q/\Z) \longrightarrow \Hom(\Z, \Q/\Z) \longrightarrow 0
\end{equation}
Let us analyze the terms in \eqref{seqB}:
\begin{itemize}[leftmargin=*, labelindent=0pt]
    \item[] $\Hom(\Z, \Q/\Z) \cong \Q/\Z$ (determined by the image of 1).
    \item[] $\Hom(\Q/\Z, \Q/\Z) \cong \Zhat$. This is the ring of endomorphisms of the torsion group $\Q/\Z$. An element $u \in \Zhat$ acts by multiplication.
\end{itemize}
Thus, \eqref{seqB} becomes:
\begin{equation}\label{seqB1}
0 \longrightarrow \Zhat \longrightarrow \Hom(\Q, \Q/\Z) \longrightarrow \Q/\Z \longrightarrow 0
\end{equation}


The map $\psi$ in \eqref{quotad1} induces a morphism between \eqref{seqA} and \eqref{seqB1}:
\[
\begin{tikzcd}
0 \arrow[r] & \Zhat \arrow[r] \arrow[d, "\psi|_{\Zhat}"] & \A_f \arrow[r, "\pi"] \arrow[d, "\psi"] & \Q/\Z \arrow[r] \arrow[d, "\text{id}"] & 0 \\
0 \arrow[r] & \Hom(\Q/\Z, \Q/\Z) \arrow[r] & \Hom(\Q, \Q/\Z) \arrow[r] & \Hom(\Z, \Q/\Z) \arrow[r] & 0
\end{tikzcd}
\]
The right vertical arrow is the identity (canonical isomorphism): for $a \in \A_f$, restricting $\psi(a)$ to $\Z$ gives $n \mapsto \pi(na)$, which is determined by $\pi(a) \in \Q/\Z$.
 
 The left vertical arrow maps $u \in \Zhat$ to the homomorphism $x \mapsto \pi(ux)$ on $\Q/\Z$. This is the canonical identification of $\Zhat$ with $\operatorname{End}(\Q/\Z)$.

By the Short Five Lemma, since the outer vertical maps are isomorphisms, the middle map $\psi$ is an isomorphism.
Next, we consider the investigation of $\psi$ in topological terms.\vspace{.02in}

\emph{Topology.}
$\A_f$ carries the restricted product topology, making it a locally compact group, while 
$\Hom(\Q, \Q/\Z)$ carries the compact-open topology. Since $\Q$ is discrete, this is the topology of pointwise convergence.
The isomorphism $\psi$ is a homeomorphism because it matches the compact open subgroups, namely:\vspace{.02in}

- The maximal compact subgroup of $\A_f$ is $\Zhat$.

- The maximal compact subgroup of $\Hom(\Q, \Q/\Z)$ is $\Hom(\Q/\Z, \Q/\Z) \cong \Zhat$ (those maps vanishing on $\Z$).

The filtration of $\A_f$ by $n^{-1}\Zhat$ corresponds exactly to the filtration of homomorphisms by the order of their restriction to $\Z$.
\end{proof}

\end{document}
